\chapter*{Abstract}
\addcontentsline{toc}{chapter}{Abstract}

Declarative process models such as Dynamic Condition Response (DCR)
graphs specify the constraints an execution must satisfy rather than a
fixed control flow, allowing flexible, compliant behaviour. This
flexibility comes at a cost: the resulting \emph{flat conformance
space} treats every compliant trace as equally valid, leaving no
mechanism to identify which executions are operationally
preferable in terms of cost, duration, or resource usage. Existing
solutions to this problem are either exact but unscalable (constraint
and satisfiability-based optimisers) or scalable but uncompliant
(standard reinforcement learning, which optimises reward without
regard for the process specification).

This thesis proposes a compliance-aware, multi-objective reinforcement
learning framework that closes this gap. A Proximal Policy Optimisation
(PPO) agent is trained inside a DCR execution engine whose shield
intercepts any action that is not currently enabled, leaving the
process marking unchanged; every trace generated during training is
therefore valid by construction, independently of reward design or
training budget. A normalised reward scalarises cost and duration under
configurable weights $(\alpha,\beta)$, so that training independent
agents across a small weight grid and pooling their accepting traces
approximates the Pareto front of compliant executions. The framework is
further extended to a role-conditioned action space, in which each
event can be executed by one of several resource roles with
differentiated cost--duration profiles, and to a capacity-constrained
variant that couples per-event role choices into a genuine
opportunity-cost allocation problem.

The framework is evaluated on two benchmark processes and externally
validated on a suite of independently mined DCR graphs never used
during development. On the Expense Report benchmark it recovers the
exact Pareto front identified by a CSP+COP exact solver; on the larger
Loan Application benchmark it discovers a well-populated, seed-stable
front of 26 non-dominated points while maintaining full compliance
throughout training, and the role-conditioned extension shows the agent
learning event-level resource specialisation consistent with the
underlying cost--duration economics. Across all experiments, structural
compliance never depends on convergence: invalid transitions are
structurally impossible from the first training episode. These results
indicate that compliance-by-construction and scalable multi-objective
optimisation are not competing requirements but can be jointly achieved
within a single reinforcement learning pipeline.
